\subsection{Slack Variables and the Slack-Relaxed LP Subproblem}
\label{sec:slacks}

\subsubsection*{Theory}

A slack variable is a non-negative auxiliary variable added to a constraint
to absorb the amount by which the constraint cannot be satisfied. Slacks
have a long history in linear programming and in the elastic-mode and
$\ell_1$-penalty SQP literature
\citep[Ch.~17--18]{NocedalWright2006},\citep{Fletcher1987}. The classical
$\ell_1$ formulation introduces a vector $s \in \mathbb{R}^m$ with
$s \geq 0$ such that
\begin{equation}
    c_k + J_{y,k}\,\Delta y + s \;\geq\; 0,
\end{equation}
and adds the linear penalty $M \sum_i s_i$ to the LP objective with
$M > 0$. The resulting slack-relaxed LP solved at iteration $k$ is
\begin{align}
    \min_{\Delta y,\, s} \quad
        & g_{y,k}^{\!\top} \Delta y \;+\; M \sum_{i=1}^{m} s_i, \\
    \text{s.t.}\quad
        & c_k + J_{y,k}\,\Delta y + s \;\geq\; 0, \\
        & s_i \geq 0, \qquad i = 1, \dots, m, \\
        & \max\!\bigl(y^L - y_k,\, -\delta_k\bigr) \leq \Delta y \leq
          \min\!\bigl(y^U - y_k,\, +\delta_k\bigr).
\end{align}
The standard upper-bound form expected by linear-programming solvers is
recovered by writing the linearized constraint as
$-J_{y,k}\,\Delta y - s \leq c_k$.

\subsubsection*{Why slacks are good}

\begin{itemize}[nosep,leftmargin=1.5em]
    \item \textbf{Always feasible LP.} The relaxed LP has at least the
          trivial point $\Delta y = 0$, $s_i = \max(0, -c_{k,i})$ in its
          feasible set. The optimizer therefore never breaks down because
          of a momentarily infeasible iterate, which is essential in
          structural problems where the starting design is rarely
          strictly feasible.
    \item \textbf{Slacks reveal information.} A large optimal slack
          $s_i^{\star}$ is a quantitative signal that constraint $i$
          cannot be satisfied locally within the present move limits.
          The solver records $\sum_i s_i$ and $\max_i s_i$ as
          diagnostics; in the implementation, persistently large slacks
          drive the penalty $M$ upwards
          (Section~\ref{sec:adaptive_penalty}), while the move limits
          respond independently to the predicted-vs-actual merit ratio
          (Section~\ref{sec:adaptive_move}).
    \item \textbf{Penalty discourages artificial feasibility.} The term
          $M \sum_i s_i$ in the LP objective makes the optimizer prefer
          steps that satisfy the linearized constraints \emph{without}
          using slack whenever possible, while still allowing slack to
          be used when the alternative is no step at all. This is the
          $\ell_1$-penalty principle of \citet{Fletcher1987} adapted to
          the LP setting.
    \item \textbf{Connection to the merit function.} The same penalty
          $M$ that weights the slack cost in the LP also weights the
          violation term in the nonlinear merit function defined in
          Section~\ref{sec:merit}. This consistency guarantees that the
          subproblem and the acceptance test agree on how much
          infeasibility is tolerable at the current iterate.
\end{itemize}

\subsubsection*{Implementation}

The slack-relaxed LP is solved with \texttt{scipy.optimize.linprog}
using the HiGHS backend \citep{Huangfu2018HiGHS}. Both the scaled
step and the slacks are stacked into a single decision vector
\begin{equation}
    z \;=\; \begin{bmatrix} \Delta y \\ s \end{bmatrix}
    \;\in\; \mathbb{R}^{n+m},
\end{equation}
so the LP is presented to the solver in the standard form
\begin{equation}
    \min_{z} \; c_{\mathrm{lp}}^{\!\top} z
    \quad \text{s.t.} \quad
    A_{\mathrm{ub}}\, z \leq b_{\mathrm{ub}},
    \qquad z_{\mathrm{lb}} \leq z \leq z_{\mathrm{ub}},
\end{equation}
with
\begin{equation}
\label{eq:lp_data}
    c_{\mathrm{lp}} =
        \begin{bmatrix} g_{y,k} \\ M\,\mathbf{1}_m \end{bmatrix},
    \qquad
    A_{\mathrm{ub}} =
        \begin{bmatrix} -\,J_{y,k} & -\,I_m \end{bmatrix},
    \qquad
    b_{\mathrm{ub}} = c_k .
\end{equation}
The block $A_{\mathrm{ub}} z \leq b_{\mathrm{ub}}$ expands to
$-J_{y,k}\,\Delta y - s \leq c_k$, which is exactly the linearized
constraint $c_k + J_{y,k}\,\Delta y + s \geq 0$. The variable box
$z_{\mathrm{lb}} \leq z \leq z_{\mathrm{ub}}$ encodes both the trust
region and the lower bound on the slacks:
\begin{align}
    \max\!\bigl(y^L - y_k,\, -\delta_k\bigr) \;\leq\; \Delta y \;&\leq\;
    \min\!\bigl(y^U - y_k,\, +\delta_k\bigr), \\
    s_i \;&\geq\; 0,
    \qquad i = 1, \dots, m .
\end{align}
The corresponding code excerpt assembles the data exactly in this layout:
\begin{quote}
\begin{verbatim}
c_lp  = np.concatenate((g_y, np.full(m, penalty)))
A_ub  = np.hstack((-J_y, -np.eye(m)))
b_ub  = c.copy()

dy_bounds = [(max(yL[i] - y[i], -delta_y[i]),
              min(yU[i] - y[i],  delta_y[i])) for i in range(n)]
bounds    = dy_bounds + [(0.0, None)] * m

res = linprog(c_lp, A_ub=A_ub, b_ub=b_ub, bounds=bounds,
              method="highs")
\end{verbatim}
\end{quote}
The first $n$ entries of \texttt{res.x} are returned as $\Delta y$ and
the trailing $m$ entries as the optimal slacks $s^\star$. Two scalar
diagnostics are stored in the iteration history,
\begin{equation}
    \mathrm{slack\_sum}_k \;=\; \sum_{i=1}^{m} s_i^\star,
    \qquad
    \mathrm{slack\_max}_k \;=\; \max_{i} s_i^\star,
\end{equation}
which feed the penalty-update rule
(Section~\ref{sec:adaptive_penalty}) and the convergence test
(Section~\ref{sec:convergence}).

\paragraph{Initial penalty.}
The penalty parameter is initialised at
$M_0 = 10^3$ (\texttt{infeasibility\_penalty}). This value is large
enough to discourage gratuitous use of slack on a well-scaled
problem, but small enough that an early-iteration LP which genuinely
cannot satisfy its linearized constraints can absorb the infeasibility
into $s$ without becoming numerically ill-conditioned. Subsequent
iterations grow $M$ geometrically by a factor of $10$
(\texttt{penalty\_increase}) up to a ceiling of $10^9$
(\texttt{penalty\_max}); the schedule is described in
Section~\ref{sec:adaptive_penalty}.

\paragraph{The unconstrained case.}
If $m = 0$ the slack block is empty: $c_{\mathrm{lp}} = g_{y,k}$, the
inequality matrix $A_{\mathrm{ub}}$ is omitted entirely, and the LP
collapses to the linearized objective minimised over the trust-region
box alone. In this case the LP is trivially feasible, since the box
always contains $\Delta y = 0$.

\paragraph{LP failure handling.}
By construction the slack-relaxed LP admits the feasible point
$\Delta y = 0$, $s_i = \max(0, -c_{k,i})$, so true infeasibility is
ruled out. In practice, however, \texttt{linprog} can still return
\texttt{success = False} for numerical reasons --- typically a trust
region so small that the LP becomes degenerate, or near-parallel
constraint rows produced by an ill-conditioned linearization. The
implementation reacts conservatively: the scaled move limits are
halved, $\delta_{k+1} = \tfrac{1}{2}\,\delta_k$, the current iterate
$x_k$ is retained, and the iteration loop continues. If the resulting
move limits fall below \texttt{move\_limit\_min} ($10^{-5}$), the
optimizer terminates with the diagnostic message ``LP failed and
scaled move limits are below \texttt{move\_limit\_min}''.
